% !TEX root = ../notes.tex

\documentclass[../notes.tex]{subfiles}

\begin{document}

\section{April 13}
Today, we continue discussing the Brauer--Manin obstruction.

\subsection{Iskovskikh's Conic Bundle}
This subsection will be interested in the following conic bundle.
\begin{example}[Iskovskikh]
	Consider the subvariety $U\subseteq\AA_\QQ^3$ cut out by
	\[U\colon y^2+z^2=\left(3-x^2\right)\left(x^2-2\right).\]
	We will show that $U$ admits no rational points using the Brauer--Manin obstruction.
\end{example}
\begin{remark}
	Iskovskikh's original proof that $U$ admits no rational points merely used quadratic reciprocity. This is not so surprising because the Brauer--Manin obstruction amounts to quadratic reciprocity in this situation. Namely, the fundamental exact sequence over $\QQ$ (for quaternion algebras) can be checked to be equivalent to quadratic reciprocity.
\end{remark}
Our first task is to turn $U$ into a nice variety, so we need to find a smooth projective model. For this, we should explain why we keep calling $U$ a conic bundle.
\begin{definition}[conic]
	A \textit{conic} over a field $k$ is a variety cut out by $s=0$ in $\PP^2_k$, where $s$ is some homogeneous polynomial in $k[x_0,x_1,x_2]$ of degree $2$.
\end{definition}
\begin{remark}
	For motivation, we give a coordinate-free definition: for a vector space $E$ over $k$, we may view $\PP^2_k=\PP E$ as hyperplanes in $E$, and we note that homogeneous polynomials $s$ of degree $2$ are elements in $\op{Sym}^2E$.
\end{remark}
\begin{definition}[conic bundle]
	A \textit{conic bundle} over a scheme $B$ is a scheme $C\to B$ which is cut out in $\PP\mc E$, where $\mc E$ is a vector bundle of rank $3$, by a nowhere vanishing section $s\in\Gamma\left(B;\op{Sym}^2\mc E\right)$.
\end{definition}
\begin{remark}
	We are not requiring our conic bundles to have any smoothness, so the fibers could easily fail to be irreducible (e.g., by being the union of two lines).
\end{remark}
\begin{example}
	If $\mc E$ is the direct sum of three line bundles $\mc L_0\oplus\mc L_1\oplus\mc L_2$, then choosing nowhere-vanishing sections $s_i\in\Gamma\left(B;\mc L_i^{\otimes2}\right)$ for each $i$ allows us to define a ``diagonal'' conic bundle with the section $s\coloneqq s_0+s_1+s_2$.
\end{example}
\begin{example}[Ch\^atelet surface] \label{ex:chat}
	Fix a field $k$ of characteristic not two, and choose a unit $a\in k^\times$ and a separable quartic $P(x)\in k[x]$ with homogenization $\widetilde P(x,w)\in k[x,w]$. Then we can define the Ch\^atelet surface $U\colon y^2-az^2=P(x)$ in $\AA^3$, which completes to a conic bundle $X\to\PP^1$ as follows: let our line bundle be $\mc E\coloneqq\OO\oplus\OO\oplus\OO(2)$, and choose our section to be $s\coloneqq(1,-a,-\widetilde P)$, which then cut out a diagonal conic bundle. Then $X$ contains $U$ as a dense open subscheme.
\end{example}
We now turn to our surface.
\begin{definition}[Iskovskikh's surface]
	We define \textit{Iskovskikh's surface} to be the Ch\^atelet surface $X\to\PP^1$ (over $\QQ$) which contains
	\[U\colon y^2+z^2=\left(3-x^2\right)\left(x^2-2\right)\]
	as a dense open subscheme.
\end{definition}
\begin{remark} \label{rem:isok-infinity}
	The fiber at infinity is cut out by $y^2+z^2=\left(3x^2-1\right)\left(1-2x^2\right)$; namely, $x=\infty$ is $y^2+z^2=-w^2$.
\end{remark}
Let's start by showing that $X$ admits local points.
\begin{proposition} \label{prop:isok-hasse}
	Iskovskikh's surface $X$ over $\QQ$ has a $\QQ_v$-point for each place $v$.
\end{proposition}
\begin{proof}
	A general approach would be to use Lang--Weil estimates to find a smooth $\FF_p$-points for large primes $p$,\footnote{Explicitly, the Lang--Weil estimates show that geometrically irreducible varieties $X$ over $\FF_q$ of dimension $d$ have $\#X(\FF_q)=q^d+O_X(q^{d-1/2})$, where the constant in $O_X$ depends explicitly on the geometry of $X$.} and then we use Hensel's lemma to lift these $\FF_p$-points to $\QQ_p$-points. Once we have handled large primes, we can explicitly find points for small $p$, and we can handle the archimedean place separately.

	Here is a more general argument: we can show that the subset
	\[U\colon y^2+z^2=\left(3-x^2\right)\left(x^2-2\right)\]
	already admits $\QQ_v$-points. We now have the following cases.
	\begin{itemize}
		\item For finite primes $p\ge3$, we may set $x\coloneqq1$ and then hunt for $(y,z)\in\ZZ_v^2$ for which $y^2+z^2=-2$. It is enough to find a smooth $\FF_p$-point, but by inspection, we see that all $\FF_p$-points will be smooth (because $y$ and $z$ cannot simultaneously vanish). For this, we note that $y^2$ takes $\frac{p+1}s$ values as $y\in\FF_p$ varies, and the same is true for $z^2$, so there must be some intersection between the values of $y^2$ and $-2-z^2$ as $y$ and $z$ varies.
		\item For $p=2$, we have the point $(0,1,\sqrt{-7})$
		\item For the real place, we have the point $(\sqrt2,0,0)$.
		\qedhere
	\end{itemize}
\end{proof}
We now begin the more difficult task of showing that $X$ admits no rational point. To start, we need to find an interesting class in $\op{Br}X$.
\begin{lemma} \label{lem:isok-quaternion}
	Let $X$ be Iskovskikh's surface over $k\coloneqq\QQ$, and set $K\coloneqq k(X)$. Then the quaternion algebra $A\coloneqq\left(3-x^2,-1\right)\in\op{Br}K$ is in $\op{Br}X$.
\end{lemma}
\begin{proof}
	The idea is to rewrite $A$ in a way that gives it a definition everywhere in $X$. Define $B\coloneqq\left(x^2-2,-1\right)$ and $C\coloneqq\left(3/x^2-1,-1\right)$. Then $A$, $B$, and $C$ are all two-torsion elements (because they are quaternion algebras). Some arithmetic with quaternion algebras (using the fact that the sum in the Brauer group is given by the cup product of \Cref{ex:quaternion-by-cup}) then lets us write
	\[A+B=\left(\left(3-x^2\right)\left(2-x^2\right),-1\right)=\left(y^2+z^2,-1\right),\]
	but $y^2+z^2=-1$ is a norm from $K(\sqrt{-1})/K$, so this quaternion algebra vanishes. Similarly,
	\[A-C=\left(x^2,-1\right),\]
	which also vanishes for the same reason.

	We now complete the proof.
	\begin{itemize}
		\item Note that $A$ is automatically defined over $X\setminus(X_\infty\cup X_{3-x^2})$, where $X_\infty$ is the fiber over $\infty\in\PP^1$ and $X_{3-x^2}$ is the fiber over the closed point cut out by $3-x^2\in\AA^1$.
		\item Similarly, $B$ is defined over $X\setminus(X_\infty\cup X_{x^2-2})$, where $X_{x^2-2}$ is as above.
		\item Lastly, $C$ is defined over $X\setminus(X_0\cup X_{3-x^2})$.
	\end{itemize}
	The above subsets of definition cover $X$, so $A\in\op{Br}K$ has no residues anywhere on $X$, so it extends to $\op{Br}X$ by \Cref{cor:extend-brauer}.
\end{proof}
We now want to compute $X(\AA_k)^A$.
\begin{lemma} \label{lem:isok-invs}
	Let $X$ be Iskovskikh's surface over $\QQ$. For each place $v$, we compute the evaluation maps
	\[X(\QQ_v)\stackrel A\to\op{Br}\QQ_v\into\QQ/\ZZ\]
	as
	\[\op{inv}_vA(P)=\begin{cases}
		0 & \text{if }v\ne2, \\
		1/2 & \text{if }v=2.
	\end{cases}\]
\end{lemma}
\begin{proof}
	We are just testing if the quaternion algebra $A(P)_v$ is a trivial algebra over $\QQ_v$. Using \Cref{ex:quaternion-by-cup}, one can see that the quaternion algebra $\left(3-x^2,-1\right)$ is trivial at a point $P=(x,y,z)$ in $U$ if and only if $3-x^2$ is a norm in $\im\op N_{\QQ(\sqrt{-1})/\QQ}$. Of course, presenting $A$ as $\left(3-x^2,-1\right)$ does not work for all points $P$, but it does work for $P\in U$, which will be enough for now.
	\begin{itemize}
		\item Suppose $p\notin\{2,\infty\}$, and choose a point $P\in X$ with $x$-coordinate $x$. We have two cases.
		\begin{itemize}
			\item If $x=\infty$ or $v(x)<0$, then $3/x^2-1\in\ZZ_p^\times$.
			\item If $v(x)\ge0$, then $3-x^2$ or $x^2-2$ is in $\ZZ_p^\times$ (because these two elements sum to $1$).
		\end{itemize}
		It follows that $A(P)_v$ can be presented to live in $\op{Br}\ZZ_p$ (in all cases!), which vanishes for finite odd primes $p$, so $\op{inv}_pA(P)=0$.
		\item Suppose $p=\infty$. By \Cref{rem:isok-infinity}, our point $P$ cannot live over $\infty\in\PP^1$ because the conic $y^2+z^2=-w^2$ has no real points. Thus, $P\in U$ has some finite $x$-coordinate $x$. Then $3-x^2$ and $x^2-2$ sum to $1$, so at least one is positive and hence a square in $\RR^\times$.\footnote{This is now well-known, called the ``Etingof principle.''} Thus, $\op{inv}_\infty A(P)=0$.
		\item Suppose $p=2$. We have the following cases.
		\begin{itemize}
			\item If $v(x)>0$, then $3-x^2\equiv-1\pmod4$.
			\item If $v(x)=0$, then $x^2-2\equiv-1\pmod4$.
			\item Lastly, if  $v(x)>0$, then $3/x^2-1\equiv-1\pmod4$.
		\end{itemize}
		Now, any element which is $-1\pmod4$ is not a norm in the image of $\QQ_2(\sqrt{-1})\to\QQ_2$; explicitly, we are saying that $-1\pmod4$ cannot be written as a sum of two squares, which can be checked$\pmod4$. Thus, the above three cases have shown that $\op{inv}_2A(P)=1/2$ for all points $P$.
		\qedhere
	\end{itemize}
\end{proof}
\begin{remark}
	Morally, what is going on here is that the place $v$ where $X$ has good reduction automatically kill $A$ because we can then go define $A$ over $\ZZ_p$. It is only at the place $p=2$ where we have a chance of finding some nontrivial invariants.
\end{remark}
The above lemma now feeds into the following result.
\begin{theorem}
	Let $X$ be Iskovskikh's surface over $\QQ$. Then $X(\AA_\QQ)$ is nonempty, but $X(\AA_\QQ)^{\mathrm{Br}}$ is empty.
\end{theorem}
\begin{proof}
	\Cref{prop:isok-hasse} shows that $X$ has $\QQ_v$-points for every place $v$, so $X(\AA_\QQ)\ne\emp$ follows by the properness of $X$. However, using the quaternion algebra $A$ from \Cref{lem:isok-quaternion}, we find that $X(\AA_\QQ)^A$ is empty: indeed, \Cref{lem:isok-invs} implies that
	\[\sum_v\op{inv}_vA(P)=\frac12\]
	for all $P\in X(\AA_\QQ)$.
\end{proof}
\begin{remark}
	It remains a question where the magical quaternion algebra $A$ comes from. The idea is that $3-x^2$ and $x^2-2$ sum to $1$, so they are relatively prime. However, they sum to a norm $y^2+z^2$, so one expects both of them to individually be norms. Thus, the quaternion algebra $A$ should be found to be trivial most of the time, with some notable problems at $2$.
\end{remark}
\begin{remark}
	It may have been reasonable to test with Brauer elements from $\op{Br}\QQ$ for our Brauer--Manin obstruction. It turns out that such elements never provide an obstruction to rational points. Instead, for Iskovskikh's surface $X$, one can use the techniques from \Cref{subsec:brauer-computations} to show that there is a split exact sequence
	\[0\to\op{Br}\QQ\to\op{Br}X\to\ZZ/2\ZZ\to0,\]
	where the $\ZZ/2\ZZ$ lifts to $A\in\op{Br}X$; thus, $A$ turns out to have been the only possibility for $X$ to have the Brauer--Manin obstruction! As a corollary of this remark, we note that the Brauer--Manin obstruction never says anything interesting if $\op{Br}k\to\op{Br}X$ is surjective.
\end{remark}
\begin{remark}
	One does not expect the Brauer--Manin obstruction to be sharp. However, if we pass to covers, then one can show that the obstruction is sharp. Explicitly, one needs to test the Brauer--Manin obstruction on all open covers $X=\bigcup_iX_i$, where the idea is that
	\[X(k)\subseteq\bigcup_iX_i(\AA_k)^{\op{Br}X_i},\]
	where $X_i(\AA_k)^{\op{Br}X_i}$ is the Brauer set for $X_i$.
\end{remark}

\subsection{The Descent Obstruction}
Our next goal is to provide an example where there are no rational points, but this is not seen even by the Brauer--Manin obstruction. To this end, we need a different way to show that there are no rational points, for which we use descent obstruction. The starting point is the following classification.
\begin{definition}[algebraic group]
	Fix a field $k$. Then an \textit{algebraic group} $G$ over $k$ is a smooth affine group scheme of finite type over $k$.
\end{definition}
\begin{theorem}
	Fix a field $k$, and let $G$ be an algebraic group. For a $k$-scheme $X$. The $G$-torsors over $X$ are parameterized by $\mathrm H^1(X;G)$.
\end{theorem}
\begin{proof}
	Omitted!
\end{proof}
We now retell our Brauer--Manin obstruction.
\begin{definition}
	Fix a nice variety $X$ over a field $k$. Fix an algebraic group $G$ and a $G$-torsor $f\colon Y\to X$. Then we define the \textit{evaluation map} $\op{ev}_Y\colon X(k)\to\mathrm H^1(k;G)$ as taking the point $x\in X(k)$ the pullback of $Y\in\mathrm H^1(X;G)$ along the map $x\colon\Spec k\to X$. Equivalently, $\op{ev}_Y(x)$ is the fiber $f^{-1}(\{x\})\to\Spec k$.
\end{definition}
\begin{remark}
	There is a decomposition
	\[X(k)=\bigsqcup_{\sigma\in\mathrm H^1(k;G)}\{x\in X(k):\op{ev}_Y(x)=\sigma\}.\]
	Because a torsor is trivial if and only if it admits a rational point, we see that the set of $x\in X(k)$ where $\op{ev}_Y(x)$ is trivial is simply $f(Y(k))$. More generally, for a class $\sigma$, there is a way to twist the torsor $f\colon Y\to X$ to the torsor $f^\sigma\colon Y^\sigma\to X$ (which is basically subtraction fiber-wise). Tracking around the subtraction, we find that
	\[X(k)=\bigsqcup_{\sigma\in\mathrm H^1(k;G)}f^\sigma(Y^\sigma(k)).\]
\end{remark}
The descent obstruction now observes that
\[X(k)\subseteq\bigsqcup_{\sigma\in\mathrm H^1(k;G)}f^\sigma(Y^\sigma(\AA_k)).\]
It turns out that $Y^\sigma(\AA_k)$ is nonempty for only finitely many $\sigma$, and one can even effectively compute which $\sigma$ we have to study.
\begin{definition}[descent set]
	Fix a nice variety $X$ over a field $k$ and an algebraic group $G$. For a $G$-torsor $f\colon Y\to X$, we define $X(\AA_k)^f$ to be the union
	\[X(\AA_k)^f\coloneqq\bigsqcup_{\sigma\in\mathrm H^1(k;G)}f^\sigma(Y^\sigma(\AA_k)).\]
	We define the \textit{descent set} $X(\AA_k)^{\mathrm{descent}}$ to be the intersection of $X(\AA_k)^f$ as $f$ varies over all $G$-torsors and $G$ varies over all smooth algebraic groups.
\end{definition}
\begin{remark}
	One can recover the weak Mordell--Weil theorem for an elliptic curve $E$ by taking $G=E[2]$.
\end{remark}

\end{document}